% Pillar-3 (group size) standalone abstract
\begin{abstract}
The group size $G$ is GRPO's most distinctive hyperparameter: it sets how many
completions are sampled per prompt and, through the group-relative baseline, how
much preference signal each prompt contributes. We study the effect of $G$ as one
pillar of \ours{}, a benchmark of 70+ RL runs across seven libraries and
five model families ($0.6$B--$\sim$671B), centered on GSM8K with HumanEval and
tool-use checks. Two results
anchor the pillar. First, trainability varies with $G$ in a non-monotone way:
intermediate group sizes often behaved better than the smallest or largest we
tested, but the data do not justify a universal optimum. Our measured equivalence
tests extend to $G{=}16$; on the near-ceiling arithmetic task, held-out accuracy
is essentially flat across this range---a retention \emph{ratio} of ${\approx}1.00$
between $G{=}2$ and $G{=}16$ (i.e.\ $G{=}16$ lands within noise of $G{=}2$; this is
a ratio, not an accuracy above $100\%$). We caution that this is a \emph{saturated}
regime: near the ceiling, variance is compressed, so the finding is ``no detectable
$G$ effect under saturation,'' not evidence of equivalence on harder, unsaturated
tasks. The effective signal-to-noise ratio
rises only sublinearly with $G$ (about $52\%$ of the $\sqrt{G}$ ideal). We do
\emph{not} report a directly trained, matched-budget $G{=}32$ cell in the
primary sweep: the paper's $G{=}32$ values come from an explicitly labeled
reconstructed token-budget grid and are hypotheses, not measurements. A separate
505-task conditional utility audit, using $(1-\mathrm{ZVF})/\sqrt{G}$ as the
declared cost proxy, selects $G{=}4$ over $G{=}5$ by $0.00177$ with a 95\%
bootstrap interval $[0.00082,0.00270]$; that result is cohort- and
objective-conditional, not a universal optimum. Second, we make explicit the sense in which GRPO
\emph{admits an all-pairs preference (DPO-like) interpretation}: the group-centered
update behaves as an implicit all-pairs preference
contrast---consistent with logged per-step diagnostics (advantage variance,
ZVF, and gradient norms), though direct gradient-vector validation is still future
work, so we frame this as an interpretation rather than a proven equivalence---with zero-variance groups
projected out. We frame $G$ as a preference-density dial rather than a variance
knob, test equivalence with TOST and difficulty-stratified retention, and release
the in-repository artifacts used here.
\end{abstract}
